\subsection{Preliminaries}

Let $\mathcal{U}$, $\mathcal{B}$, $\mathcal{L}$ and $\mathcal{V}$ be pairwise disjoint
sets of IRIs, blank nodes, literals and variables. An \emph{RDF triple} is an element of
$\Tt = (\mathcal{U} \cup \mathcal{B}) \times \mathcal{U} \times
(\mathcal{U} \cup \mathcal{B} \cup \mathcal{L})$, and an \emph{RDF graph} a finite subset
of $\Tt$. A \emph{triple pattern} is an element of
$\mathcal{P} = (\mathcal{U} \cup \mathcal{V}) \times (\mathcal{U} \cup \mathcal{V})
\times (\mathcal{U} \cup \mathcal{L} \cup \mathcal{V})$ and a \emph{basic graph pattern}
(BGP) a finite set of triple patterns.

A \emph{solution mapping} is a partial function
$\mu : \mathcal{V} \to \mathcal{U} \cup \mathcal{B} \cup \mathcal{L}$; we write
$\mu_\emptyset$ for the one with empty domain, $\Omega$ for a set of solution mappings,
and $\bowtie$ for their join. Following \citet{perez2009semantics}, for a triple pattern
$t$ we write $\mu(t)$ for the triple obtained by replacing the variables in $t$ according
to $\mu$, and $\mathrm{var}(t)$ for the variables occurring in $t$. A triple pattern
evaluates over a graph $G$ to
\[
  \ev{t}{G} \;=\; \{\, \mu \mid \mathrm{dom}(\mu) = \mathrm{var}(t) \text{ and }
  \mu(t) \in G \,\},
\]
and evaluating a BGP $B$ distributes over its patterns,
\[
  \ev{B}{G} \;=\; \bigjoin_{t \in B} \ev{t}{G},
  \qquad \ev{\emptyset}{G} = \{\mu_\emptyset\},
\]
so the empty BGP evaluates to the identity for $\bowtie$. For a BGP $B$ we write
\[
  \mu(B) \;=\; \{\, \mu(t) \mid t \in B \,\}
\]
for the set of triples obtained by applying $\mu$ to every pattern of $B$; patterns that
$\mu$ maps to the same triple contribute it once. 

As triple patterns contain no blank nodes, this evaluation coincides with basic graph
pattern matching in SPARQL~\cite{sparql11-query}. Writing
$\mathrm{var}(B)$ for the variables of $B$, a mapping is a solution of $B$ exactly when
it binds all variables of $B$ and every triple it produces from $B$ occurs in $G$,
\begin{equation}
  \label{eq:bgp-match}
  \mu \in \ev{B}{G} \iff \mathrm{dom}(\mu) = \mathrm{var}(B) \text{ and }
  \mu(B) \subseteq G.
\end{equation}
In other words, $\mu(B)$ is the set of triples of $G$ from which the solution $\mu$ is
obtained.

Blank nodes are existential variables scoped to the document that serialises
them~\cite{rdf11concepts}: no two sources share one, and none can be referenced in a
later request~\cite{prudhommeaux2013federated}. As a result, a query decomposed over
several sources can lose answers that join on a blank node within one
source~\cite{stolpe2017blank}. We exclude them.

\begin{assumption}[Skolemised sources]
\label{as:skolem}
Every source publishes a graph free of blank nodes,
$\dataf(d) \subseteq \mathcal{U} \times \mathcal{U} \times (\mathcal{U} \cup \mathcal{L})$.
Blank nodes are replaced by Skolem IRIs~\cite{rdf11concepts}, and a source republishing
data preserves those minted by its origin.
\end{assumption}

FedQPL assumes only that members use pairwise disjoint blank
nodes~\cite{cheng2020fedqpl}. This is exactly what breaks aggregation: a composite
republishing a document must rename its blank nodes apart, so they no longer denote
the nodes its origin published. Skolem IRIs survive republication.

